% !TEX program = xelatex
% Jeffrey Alan Scudder — Selected Work (Serpentine FAE Fellowship portfolio)
% Build:  tectonic -X compile portfolio.tex   (or xelatex portfolio.tex)
% Matches the AC paper / CV house style: Latin Modern + AC palette.

\documentclass[11pt,letterpaper]{article}

\usepackage[top=0.7in, bottom=0.7in, left=0.9in, right=0.9in]{geometry}

% === FONTS (Latin Modern via bundled OTFs — works under tectonic & MacTeX) ===
\usepackage{fontspec}
\setmainfont[
  Extension = .otf,
  UprightFont = lmroman10-regular,
  BoldFont = lmroman10-bold,
  ItalicFont = lmroman10-italic,
  BoldItalicFont = lmroman10-bolditalic,
]{lmroman10}
\setsansfont[
  Extension = .otf,
  UprightFont = lmsans10-regular,
  BoldFont = lmsans10-bold,
  ItalicFont = lmsans10-oblique,
]{lmsans10}
\setmonofont[
  Extension = .otf,
  UprightFont = lmmono10-regular,
  ItalicFont = lmmono10-italic,
  Scale = 0.86,
]{lmmono10}

\usepackage{xcolor}
\usepackage{graphicx}
\usepackage{titlesec}
\usepackage{enumitem}
\usepackage{hyperref}
\usepackage{microtype}
\usepackage{parskip}
\usepackage{adjustbox}

% === AC palette (shared with arxiv-* papers + cv) ===
\definecolor{acpink}{RGB}{180,72,135}
\definecolor{acpurple}{RGB}{120,80,180}
\definecolor{acdark}{RGB}{64,56,74}
\definecolor{acgray}{RGB}{119,119,119}
\definecolor{acrule}{RGB}{210,205,214}

\hypersetup{colorlinks=true, linkcolor=acpurple, urlcolor=acpurple}
\setlength{\parindent}{0pt}
\pagestyle{empty}

% A single work entry: image, then title/meta/caption block.
% #1 image  #2 title  #3 year  #4 role  #5 link-text  #6 link-url  #7 caption
\newcommand{\work}[7]{%
  \begin{adjustbox}{center,cfbox=acrule 0.8pt 0pt}
    \includegraphics[width=\textwidth, height=0.52\textheight, keepaspectratio]{#1}%
  \end{adjustbox}\par
  \vspace{0.9em}
  {\sffamily\bfseries\Large\color{acdark} #2}\quad{\sffamily\color{acgray} #3}\par
  \vspace{0.25em}
  {\sffamily\small\color{acpink} #4}\hfill{\ttfamily\small\href{#6}{#5}}\par
  \vspace{0.5em}
  \begin{minipage}{\textwidth}\color{acdark}#7\end{minipage}\par
}

\begin{document}

% ===================== COVER =====================
\thispagestyle{empty}
\vspace*{0.10\textheight}
\begin{center}
  {\sffamily\color{acgray}\large selected work}\\[0.5em]
  {\bfseries\fontsize{34}{38}\selectfont\color{acdark} Jeffrey Alan Scudder}\\[0.7em]
  {\sffamily\large\color{acpurple} Aesthetic Computer — a creative computing commons}\\[1.4em]
  {\color{acgray}\rule{0.42\textwidth}{0.6pt}}\\[1.2em]
  \begin{minipage}{0.80\textwidth}\centering\color{acdark}
    A runtime, an operating system, and a social network for creative
    computing, built and given away over six years. People make pieces in
    JavaScript or KidLisp, alongside paintings, tapes, and chats — each living
    at a web address anyone can open — and a surplus laptop boots into the whole
    thing in about seven seconds. The five works that follow trace the commons
    from a single piece down to the bare metal, and back out to the people who
    gather around it.
  \end{minipage}\\[1.6em]
  {\sffamily\small\color{acgray}
    aesthetic.computer \quad·\quad Serpentine FAE Fellowship — Art\,$\times$\,Convergence \quad·\quad 2026}
\end{center}
\vfill
\begin{center}{\sffamily\footnotesize\color{acgray} Los Angeles, California}\end{center}
\clearpage

% ===================== 1. AESTHETIC COMPUTER / KIDLISP =====================
\work{img/ac-kidlisp.png}{Aesthetic Computer}{the system / a KidLisp piece · 2020–}{author}%
{aesthetic.computer}{https://aesthetic.computer}%
{A runtime, OS, and social network where artists make pieces in JavaScript or
KidLisp, alongside paintings, tapes, and chats --- each living at a shareable
web address. The piece shown is live KidLisp: its source, its rendered output,
and the short \texttt{\$code} that lets anyone open it are all on screen at
once. This is the cultural commons the proposal sets out to re-found and
sustain.}
\clearpage

% ===================== 2. AC NATIVE OS =====================
\work{img/ac-native.png}{AC Native OS}{re-founding the personal computer · 2025–}{author}%
{aesthetic.computer}{https://aesthetic.computer}%
{A bare-metal Linux that boots any surplus laptop into the creative instrument
in about seven seconds, with your identity living in the boot partition rather
than a cloud login. A retired fifty-dollar ThinkPad becomes a personal computer
you own outright --- and drops the cost of a networked laptop ensemble by two
orders of magnitude, putting a planetary laptop orchestra within reach.}
\clearpage

% ===================== 3. NOTEPAT =====================
\work{img/notepat.png}{notepat}{instrument · 2024}{author}%
{aesthetic.computer/notepat}{https://aesthetic.computer/notepat}%
{A musical instrument that turns a computer or phone keyboard into a playable,
shareable score --- tunes you can type, so a song is something you write down
and pass along. It is open-source, lives at a web address, and reached the
front page of Hacker News: a tool built for other people to depend on and make
their own work with.}
\clearpage

% ===================== 4. LAKLOK (community) =====================
\work{img/laklok.png}{laklok}{the social layer · 2024–}{author}%
{aesthetic.computer/laklok}{https://aesthetic.computer/laklok}%
{A live, public chat room that is itself an Aesthetic Computer piece ---
messages, images, and links shared in real time at its own web address. The
everyday social layer of the commons, where a community gathers, talks, and
passes work back and forth. The proposal's laptop orchestra grows out of rooms
exactly like this one.}
\clearpage

% ===================== 5. WHISTLEGRAPH =====================
\work{img/whistlegraph.png}{The Longest Whistlegraph Ever (so far)}{2022 · New Museum (Rhizome commission)}{co-creator, performer}%
{sites.rhizome.org/the-longest-whistlegraph-ever-so-far}{https://sites.rhizome.org/the-longest-whistlegraph-ever-so-far/}%
{Live, social drawing performed with Camille Klein and Alex Freundlich as
Whistlegraph, premiered at the New Museum. A whistlegraph is an audio-visual
work performed by hand --- drawing and singing a reproducible score into a
poetic image. The performance-and-community lineage behind building tools
people gather around and play together.}

\end{document}
